% !TeX root = proposal.tex
\section{Introduction}
\label{sec:intro}

% \textcolor{blue}{Reviewer 3: The first section is very verbose and could be more efficiently described.}

% \textcolor{blue}{(Response: We re-wrote Introduction to deliver the core message.)}

Edge and high-performance computing (HPC) are mutually beneficial components in large-scale scientific experiments. Data volume and computing requirements continue to grow, creating pressure on the computing and network capability of HPC systems. Edge computing can be introduced to process data at the source, lowering the need for sending data to an HPC production system and improving performance. As such, the ability to connect edge resources to an HPC system in a common scientific workflow enables large simulations running on HPC to take filtered data streamed from the edge, reducing latency and infrastructure cost, while still benefiting from the large computing power of HPC for science discovery. The intermediate artifacts of edge computing also provide immediate feedback to users, allowing them to control data processing or even the experiment itself as it occurs. Similarly, algorithms and artificial intelligence (AI) models running on edge computing can be retrained with new data from the experiment on connected HPC resources, adapting the workflow as new scientific needs arise.

The collaboration between edge and HPC, named edge-to-HPC, forms an integrated infrastructure and brings new challenges, ranging from HPC workflow management~\cite{dube2021future, badia2024integrating} to real-time control of sensors~\cite{kandel2023demonstration}. These challenges need attention for designing a holistic control of resources in the edge-to-HPC continuum. Application runs and data processing in edge computing are manually executed and curated by users, and their resource requirements are manually controlled and are not propagated outside the edge. We claim that the challenges in edge computing remain as yet unexplored for finding the potentials and limitations toward a fully connected continuum.

To support the transformative edge-to-HPC computing continuum, we describe here a project that investigates (1) how resources on edge devices can be arbitrated dynamically to ensure the performance for experiment-critical applications while being flexible for reconfiguration, (2) how system software can help design AI@Edge applications that leverage HPC platforms to tune machine learning (ML) model performance, and (3) how flexible the user-agnostic system software policies can be, in order to be deployed across a wide range of use cases and hardware setups.

In this paper we focus on the management of edge computing resources under a user-provided performance goal. The contributions of this paper are as follows. First, we identify an existing scientific workflow from a group of collaborating scientists that requires resource management and arbitration at the edge for science-driven control. Second, we design a framework that captures the performance of both system and user applications while the workflow is running and enables a control algorithm to maximize application throughput. Third, we demonstrate how such a control algorithm can be designed and how such capability for adaptive edge-computing performance can fit in the continuum, identifying challenges for future resource management on the scientific edge. The framework is open-sourced in Github\footnote{https://github.com/perarnau/charon}.

The following sections describe the scientific workflow that explains our motivation and opportunities for edge-to-HPC, the proposed framework for edge computing, experiment results that highlight the mechanism of the control, and a summary of our approach and future direction.